\section{Dokumentation der Prüfräder}
\label{sec:dokumentation_pruefraeder}

Jedes Prüfrad wurde vor dem Einbau in den Prüfstand vermessen. Die Messungen erfassen die Fertigungsstreuung der gespritzten Räder und den Zustand der Zahnkontur vor dem Laufversuch; daneben lassen sich zahnindividuelle Abweichungen dem Fehlerfenster zuordnen. Als Fehlerfenster werden die Zähne 36 bis 40 des Rades bezeichnet, in denen die bisherigen Laufversuche versagten. Zahn 1 ist an den Teilen markiert und wird von Messmaschine und Prüfstand gleich gezählt. Vermessen wurden 36 Räder und 8 Ritzel. Die zehn Räder der eigenen Laufversuche bilden die Kerngruppe, die übrigen 26 Räder einschließlich des Zusatzversuchs bei $12$\,Nm die Gruppe der Fertigungsstreuung. Sieben Ritzel stammen aus der Serie des Prüfstandsritzels, das Ritzel 86481 ist ein älteres Teil mit einem Konturscan aus dem Jahr 2023.

Die Verzahnungsmessung wurde auf der Klingelnberg P40 mit dem Messprogramm GINA durchgeführt. Je Flanke wurden am Rad fünf und am Ritzel drei Zähne in Profil und Flankenlinie abgetastet, Teilung und Rundlauf an allen Zähnen. Das Messprogramm gibt die Abweichungen nach DIN ISO 1328-1 \cite{DINISO1328-1} aus, am Profil $f_{\mathrm{H}\alpha}$, $f_{\mathrm{f}\alpha}$ und $F_{\alpha}$, an der Flankenlinie $f_{\mathrm{H}\beta}$, $f_{\mathrm{f}\beta}$ und $F_{\beta}$. An der Teilung werden $f_{\mathrm{p}}$ und $F_{\mathrm{p}}$ ausgegeben, dazu die Rundlaufabweichung $F_{\mathrm{r}}$ sowie der Kopf- und der Fußkreisdurchmesser. Die Zahnweite $W_{\mathrm{k}}$ und das Zweikugelmaß $M_{\mathrm{dK}}$ sind Prüfmaße der Zahndicke nach DIN 21773 \cite{DIN21773}; am Ritzel wird außerdem die Kopfrücknahme gemessen. Die Räder wurden in drei Runden vom 8. bis 13. Mai, am 27. Mai und am 10. Juni 2026 gemessen, 21 Räder in mehr als einer Runde. Aufspannung und Ablauf blieben in allen Runden gleich, der Bediener wechselte von Runde zu Runde.

Auf derselben Maschine wurden Konturscans im Stirnschnitt über drei Zähne und die Lücken dazwischen aufgenommen. Gescannt wurden die Zähne 1, 18 und 36 an 24 Rädern der Streugruppe und an zwei Ritzeln sowie Zahn 1 an vier Rädern der Kerngruppe. Sechs Räder der Kerngruppe besitzen keinen Scan. Das Ritzel 95912 wurde bei der für das ältere Ritzel festgehaltenen Maschinenhöhe von $87{,}57$\,mm gescannt. 30 Scans an zehn Rädern enthalten nur etwa 400 Punkte und werden als grobe Scans getrennt geführt.

Die Rauheit der Zahnflanken wurde mit einem Tastschnittgerät von Hommel-Etamic an beiden Flanken von 35 Rädern bestimmt, je Flanke mit drei Tastschnitten; die Ritzel wurden nicht gemessen. \Note{PH: Gerätetyp, Taststrecke und Grenzwellenlänge der Rauheitsmessung.} Ausgewertet werden $R_{\mathrm{a}}$ und $R_{\mathrm{z}}$ nach DIN EN ISO 21920-2 \cite{DINENISO21920-2}; die übrigen Kenngrößen des Gerätes werden unverändert mitgeführt. Die Messblätter der Verzahnungsmessung und die Rauheitsprotokolle stehen in \secref{sec:anhang_g}.

Von den drei gescannten Zahnköpfen des Ritzels 95912 liegen Mikroskopbilder mit einem Maßstabsbalken von $200$\,\textmu m vor, aufgenommen im Schrägblick auf Kopfzylinder und Flanke. Der Blickwinkel ist nicht festgehalten, deshalb dienen die Bilder dem Vergleich der Kopfkante mit dem Konturscan und nicht der Messung eines Radius. \Note{PH: Mikroskop, Gerät und Vergrößerung.} \tabref{tab:messumfang} stellt den Messumfang je Messart zusammen.

\begin{table}[htbp]
    \caption{Messumfang der Dokumentation der Prüfräder nach dem Inventar der Messdateien: Messart, Gerät, Messgrößen, Norm, Anzahl der Teile und Messungen.}
    \label{tab:messumfang}
    \centering
    \footnotesize
    \setlength{\tabcolsep}{4pt}
    \begin{tabular}{L{2.4cm}L{2.0cm}L{4.3cm}L{2.2cm}L{3.9cm}}
        Messart & Gerät & Messgrößen & Norm & Teile und Messungen \\
        \hline
        Verzahnungs\-messung (GINA) & Klingelnberg P40 & $f_{\mathrm{H}\alpha}$, $f_{\mathrm{f}\alpha}$, $F_{\alpha}$, $f_{\mathrm{H}\beta}$, $f_{\mathrm{f}\beta}$, $F_{\beta}$, $f_{\mathrm{p}}$, $F_{\mathrm{p}}$, $F_{\mathrm{r}}$; $W_{\mathrm{k}}$, $M_{\mathrm{dK}}$, $d_{\mathrm{a}}$, $d_{\mathrm{f}}$; Kopfrücknahme des Ritzels & DIN ISO 1328-1, DIN 21773 & 36 Räder, 7 Ritzel; 68 Messwertdateien mit Kurvendateien, davon 58 am Rad und 10 am Ritzel; 47 Messblätter; Rad drei Runden, Ritzel eine Runde \\
        Konturscan & Klingelnberg P40 & Stirnschnittkontur der Zähne 1, 18 und 36; Kopfkante, Kopfrücknahme, $s_{\mathrm{b}}$, $d_{\mathrm{a}}$, $d_{\mathrm{f}}$ & Sollkontur nach DIN ISO 21771 & 28 Räder, 2 Ritzel; 82 Scans, davon 30 grobe Scans an zehn Rädern \\
        Rauheit & Hommel-Etamic & $R_{\mathrm{a}}$, $R_{\mathrm{z}}$; weitere Kenngrößen unausgewertet & offen & 35 Räder, linke und rechte Flanke; 70 Messreihen zu je drei Tastschnitten \\
        Mikroskop & offen & Kopfkante im Schrägblick, Maßstab $200$\,\textmu m & keine & 1 Ritzel, Zähne 1, 18 und 36; 3 Bilder \\
    \end{tabular}
\end{table}

Ein eigenes Auswerteskript liest die Exportdateien der Geräte und ordnet jede Datei über Dateinamen und Dateikopf einem Teil, einer Messart und einem Zahn oder einer Flanke zu. Was die Dateien nicht sagen, wie die Gruppeneinteilung oder das Fehlerfenster, steht in einer getrennten Zuordnungsdatei. Die Messrunden werden aus den Datumsangaben der Dateiköpfe gebildet, weil die Namenszusätze der Wiederholmessungen uneinheitlich sind; Datumswerte mit höchstens fünf Tagen Abstand bilden eine Runde. Auffälligkeiten wie doppelte Exporte oder ein Scan mit zwei Blöcken werden im Inventar benannt und nicht bereinigt. Die zwölf Dateien einer ME-Nummer, die nur in Rauheitsdateien vorkommt, bleiben unzugeordnet.

Die Statistik über die Teile nimmt je Teil die letzte Messung; frühere Runden gehen in eine Tabelle der Wiederholbarkeit ein, damit die Fertigungsstreuung neben der Messstreuung steht. Die im Messprotokoll hinterlegten Toleranzen sind Einstellungen des Messprogramms, Qualität 7 für die Räder und 6 für die Ritzel, und werden als Toleranz laut Messprotokoll geführt. Eine Verzahnungsqualität nach Norm wird nicht berechnet. Die Sektorgrößen $F_{\mathrm{pk}}$ und $F_{\mathrm{pSk}}$ nach Anhang D der Norm werden aus den Teilungstabellen der Kurvendateien berechnet, weil der Programmwert Fpz/8 keiner der beiden entspricht. Aus denselben Tabellen werden die Werte der Zähne 36 bis 40 dem übrigen Umfang gegenübergestellt.

Je Konturscan werden die Achse, eine Evolvente je Flanke mit dem Grundkreis der Sollverzahnung, ein Kopfkreisradius je Zahnkopf und die Teilung angepasst. Eine robuste Fehlerfunktion hält Fußrundung und Kopfkante aus der Anpassung heraus; Sollkontur ist die STplus-Erzeugung mit mittlerem Zahndickenabmaß. Die Kopfkante wird als tangentialer Bogen und als symmetrische Fase über ein Größenraster angepasst und mit dem RMS-Wert beider Formen ausgewiesen. Beide Formen kommen in Betracht, weil die Zeichnung des Rades eine Rundung r $0{,}2$ verlangt, die STplus-Sollkontur aber eine Fase führt. Die Kopfrücknahme wird als Gerade über der Wälzlänge mit dem Betrag $C_{\mathrm{a}}$ am Kopf und dem Beginn bei $d_{\mathrm{Ca}}$ nach DIN ISO 21771 \cite{DINISO21771} gelesen. Weil die freie Anpassung eine Profil-Winkelabweichung als Achsverschiebung aufnimmt, wird die Achse entlang der Mittellinie des mittleren Zahnes verschoben, bis der Median der Kopfkreisradien dem GINA-Kopfkreis der frühesten Runde entspricht. Das Ritzel 86481 ohne GINA-Messung behält die freie Achse. Der Verschleißvergleich nach dem Laufversuch überlagert zwei Scans desselben Zahnes auf dieselbe Weise. \Note{PH: Nachher-Scans der Laufversuchsräder.}

Für die Rauheit gilt der Gerätemittelwert der drei Tastschnitte als Wert einer Flanke; er wird gegen die Einzelschnitte im Rahmen der gedruckten Nachkommastellen geprüft. Doppelte Exporte werden an gleichen Zahlenwerten erkannt und entfernt, nicht am Dateinamen.

% Herkunft:
% Absatz 1 (Zweck, Teile, Gruppen, Fehlerfenster): ADR-117 Context und Consequences (architecture_decisions.md
%   Zeilen 1949 bis 2003); parts.yaml (provenance.notes, groups, labels, failure_window, drawings);
%   inventory.md Kopfzeile (44 Teile, 688 zugeordnete Dateien, 12 nicht zugeordnete).
% Absatz 2 (GINA): ADR-117 Context (fünf Zähne Rad, drei Ritzel, P40); gina.py Docstring (Sizes WK, MdK,
%   Kopf, Fuß; fKo des Ritzels); Normbezeichnungen und Symbole aus quantities.yaml (Einträge
%   DINISO1328-1:2018, Tabelle 1 S. 9 bis 11; span_measurement W_k, M_dK nach DIN 21773); Runden und
%   Mehrfachmessungen aus inventory.md (Spalten GINA-Runden, mew/mka: 21 Räder mit 2 oder 3 Dateien,
%   58 Rad + 10 Ritzel = 68 .mew, 47 Messblatt-PDF); Bedienerwechsel: ADR-117 Amendment 2026-10-10 Punkt 4;
%   Runden-Daten: inventory.py Docstring (8 bis 13 Mai, 27 Mai, 10 Juni 2026) und inventory.md (Ritzel
%   2025-10-10/13, eine Runde, nur in der Tabelle genannt).
% Absatz 3 (Konturscan): contour_scan.py und measurements.py Docstrings (Stirnschnitt dreier Zähne);
%   Zahnnummern und Teile aus inventory.md (24 Räder z1/z18/z36, vier Kernräder z1, 86481 und 95912,
%   82 Scans, 30 grobe Scans mit Stern an zehn Rädern); Maschinenhöhe 87,57 bis 87,62 mm und Aufspannung
%   wie bei 86481: ADR-117 Amendment 2026-10-12 Punkt 3 und 4; sechs Kernräder ohne Scan: ADR-117
%   Consequences; grobe Scans standardmäßig ausgelassen: measurements.py Docstring (--include-coarse).
% Absatz 4 (Rauheit): roughness.py Docstring (Hommel-Etamic, drei Tastschnitte, fünfzehn Kenngrößen,
%   Mittelwerte, R_a/R_z, übrige roh, Protokolle Anhang); 35 Räder li/re aus inventory.md (97340 bis
%   97374, 97375 ohne); keine Ritzel: inventory.md und ADR-117 Consequences; Rohdurchlauf ohne Norm:
%   ADR-117 Decision 4 (ISO 4287 / ISO 21920 nicht im Repository).
% Absatz 5 (Mikroskop): ADR-117 Amendment 2026-10-12 Punkt 4 (F4092_95912_Z1/Z18/Z36, Maßstab 200 µm,
%   Schrägblick, Blickwinkel nicht festgehalten); inventory.md (3 JPG).
% Tabelle: Zahlen aus inventory.md (Dateiarten-Tabelle, Teiletabellen), Geräte aus ADR-117 Context,
%   Normen aus quantities.yaml und contour_scan.py Docstring (ISO 21771 für die Kopfrücknahme). Die
%   Datei-Zahlen der Rauheit (140 ASCII, 70 Protokolle, 210 Profile) enthalten die zwölf Dateien der
%   Nummer 97840, deshalb steht in der Tabelle die Reihenzahl 35 mal 2.
% Absatz 6 (Inventar, Runden, Auffälligkeiten): inventory.py Docstring; ADR-117 Decision 2, 3 und 5;
%   measurements.py Docstring (inventory, gina, contour, roughness); inventory.md Abschnitt "Nicht
%   zugeordnet" (12 Dateien 97840); anomalies.md.
% Absatz 7 (GINA-Auswertung): ADR-117 Amendment 2026-10-10 Punkt 2 bis 4; Decision 4 (Toleranz laut
%   Messprotokoll, Qualität 7/6); gina.py Docstring (Fpz/8 gegen F_pk und F_pSk); Roadmap M2 (Fenster
%   36 bis 40); Bezeichnungen F_pk, F_pSk aus quantities.yaml.
% Absatz 8 (Konturscan-Auswertung): contour_scan.py Docstring (Anpassung, robuste Fehlerfunktion,
%   Kopfrücknahme nach ISO 21771 mit C_a und d_Ca); ADR-117 Amendment 2026-10-11 Punkt 1 bis 3 (Sollzahn
%   mit mittlerem Abmaß, Bogen oder Fase, Zeichnung r 0,2 gegen STplus-Fase, Überlagerung, M5) und
%   Amendment 2026-10-12 Punkt 1, 2, 4 und 5 (Achse auf den GINA-Kopfkreis, Median, früheste Runde,
%   86481 freie Achse); Roadmap M5 (Nachher-Scans offen).
% Absatz 9 (Rauheits-Auswertung): roughness.py Docstring; ADR-117 Amendment 2026-10-11 Punkt 5.
% Zitate: ISO18 = DIN ISO 1328 (2018, "Zylinderräder - ISO-Toleranzsystem", Teilnummer im Eintrag nicht
%   genannt, vom Nutzer prüfen), DIN14 = DIN 21773 (2014), ISO14 = DIN ISO 21771 (2014); ISO21 ist
%   DIN ISO 1328 von 2021 und nicht ISO 21920, deshalb Platzhalter; ISO/TR 6336-30 nicht im Verzeichnis.
% Nomenklatur: W_k, M_dK, s_b, d_Ca, F_pk, F_pSk, f_H_alpha, f_f_alpha, f_H_beta, f_f_beta, d_a, d_f gegen
%   Nomenklatur.tex abgleichen (C_a, F_alpha, F_beta, F_p, F_r, f_p, R_a, R_z stehen bereits dort).
